\section{Alat Modern: Pengenalan Bundler dan Tooling}

\textbf{Bundler} menggabungkan banyak file JavaScript (dan CSS, gambar) menjadi satu atau beberapa file optimal untuk produksi. Ini mengatasi masalah: banyak HTTP request (lambat), dependensi antar modul, dan kompatibilitas browser. Bundler populer: \textbf{Vite} (cepat, modern, gunakan ini untuk proyek baru), \textbf{webpack} (paling matang, banyak konfigurasi), dan \textbf{Parcel} (zero-config) \cite{mdnToolsTesting}.

Ekosistem tooling modern juga mencakup: \textbf{npm}/\textbf{pnpm} (package manager), \textbf{ESLint} (linter), \textbf{Prettier} (formatter), \textbf{TypeScript} (tipe statis opsional), dan \textbf{Git} (version control). Memahami tooling ini penting untuk development workflow profesional, meskipun untuk belajar JavaScript dasar, vanilla HTML+JS sudah cukup \cite{mdnJsGuide}.

\begin{catatan}
\textbf{Vite untuk Pemula.} Jika memulai proyek baru yang memerlukan bundler, gunakan \textbf{Vite}: \code{npm create vite@latest my-app}. Vite sangat cepat (hot module replacement instan), mendukung ES modules natively, dan konfigurasi minimal. Webpack masih dominan di proyek legacy, tetapi Vite adalah pilihan standar untuk proyek baru \cite{mdnToolsTesting}.
\end{catatan}

